% ECONOMICS_PROBLEM: {"id": "F44-334", "title": "International business-cycle transmission", "jel_code": "F44", "source_id": "OP-0238"}
% !TeX program = xelatex
\documentclass[11pt,a4paper]{article}
\usepackage[margin=23mm]{geometry}
\usepackage{fontspec}
\setmainfont{Latin Modern Roman}
\usepackage{amsmath,amssymb,mathtools}
\usepackage{xcolor,enumitem,booktabs,longtable,array,xurl,hyperref}
\definecolor{muted}{HTML}{53636C}
\hypersetup{colorlinks=true,urlcolor=blue,linkcolor=blue}
\setlength{\parindent}{0pt}
\setlength{\parskip}{5pt}
\newcommand{\R}{\mathbb R}
\newcommand{\E}{\mathbb E}
\newcommand{\N}{\mathbb N}
\newcommand{\ind}{\mathbf 1}
\newcommand{\Var}{\operatorname{Var}}
\newcommand{\Cov}{\operatorname{Cov}}
\newcommand{\supp}{\operatorname{supp}}
\DeclareMathOperator*{\argmin}{arg\,min}
\DeclareMathOperator*{\argmax}{arg\,max}
\newcommand{\entrymeta}[3]{{\sffamily\small\color{muted}\textbf{\texttt{#1}}\quad|\quad #2\par #3\par}}
\newcommand{\Problemref}[1]{\texttt{#1}}
\newcommand{\JELref}[2]{\texttt{#2} (JEL \texttt{#1})}
\begin{document}
\section*{International business-cycle transmission}
\textbf{Problem ID:} \texttt{F44-334}\quad \textbf{JEL:} \texttt{F44}\par
% BEGIN ECONOMICS STATEMENT
\label{problem:OP-0238}
\label{audited:ECON-083}
{\sffamily\small\color{muted}Original subject group: Business cycles and macro dynamics\par}
\label{definitions:problem:30007519}
\textbf{Audit classification:} Background; proposed extension. The reference models inflation spillovers; the output-correlation network intervention is a new specification. Verification concerns the cited version, not first priority or proof that this exact editorial specification remains unsolved on October 8, 2026.
\subsection*{Definitions and precise research target}
\paragraph{Editorial mathematical specification.} Fix countries \(c=1,\ldots,C\) with measured sectoral input–output links, bilateral trade costs, cross-border assets, and country-specific policy rules. Let \(\Omega\) denote the global production-input matrix, partitioned into domestic and international blocks. A structural model \(M\) maps sectoral and financial innovations into country log outputs \(y_c\). Construct \(\Omega^0\) by replacing imported intermediate inputs with specified domestic substitutes, holding initial final-demand shares and aggregate financial shocks fixed; all resulting equilibrium prices and quantities must adjust within the model.
For horizon \(H\), define the network contribution to cycle synchronisation between countries \(c\) and \(d\) as
\[
\Delta_{cd}(H)=\operatorname{Corr}_M(\Delta_H y_c,\Delta_H y_d;\Omega)
-\operatorname{Corr}_M(\Delta_H y_c,\Delta_H y_d;\Omega^0),
\]
where \(\Delta_H y_c=y_{c,t+H}-y_{c,t}\), both correlations use the same declared innovation distribution, and each output change has strictly positive finite variance under both networks. Identify \(\Delta_{cd}(H)\), or a sensitivity set, using shocks with independently supported geographic incidence. Distinguish propagation from common shocks and endogenous network formation; specify whether national policy rules remain fixed or optimally respond. Validate predicted transmission against held-out trade disruptions. Removing international links is a model-based intervention whose substitute technology must be disclosed; its correlation effect cannot be read directly from observed trade–cycle correlations.

The matrix \(\Omega\) must refer either to fixed technology coefficients or to observed benchmark expenditure shares; in the latter case the model updates shares after intervention. A valid replacement \(\Omega^0\) specifies domestic suppliers, resource use, substitution possibilities, and trade costs, so that removed foreign inputs are not replaced for free. At a fixed initial law define \(U_c^v=y_{c,t+H}^v-y_{c,t}^v\), \(v\in\{0,1\}\), using common primitive shocks. Each correlation is \(\operatorname{Cov}(U_c^v,U_d^v)/(\operatorname{Var}(U_c^v)\operatorname{Var}(U_d^v))^{1/2}\). The identified set consists of \(\Delta_{cd}(H)\) over models matching the same observable trade, price, and output law and restrictions. This identifies a chosen network intervention, not a source-independent percentage of international synchronization.
\subsection*{Variants and assumption changes}
The following are \emph{editorial variants}, not additional published open conjectures.
\begin{enumerate}
\item \emph{Shock propagation.} Replace correlations by a foreign-output impulse response to a valid country-specific productivity shock. Hold the shock variance fixed and compare domestic-only versus international intermediate-input links; this avoids contamination from varying common-shock variance.
\item \emph{Financial links.} Retain trade links and instead remove cross-border credit using a specified domestic financing replacement. Compare trade-only, finance-only, and joint interventions, allocating their interaction explicitly.
\end{enumerate}
\subsection*{References and source audit}
\begin{itemize}
\item Julian di Giovanni; \c{S}ebnem Kalemli-\"{O}zcan; Alvaro Silva; Muhammed A. Y\i{}ld\i{}r\i{}m. \emph{Pandemic-Era Inflation Drivers and Global Spillovers}. 2023. Locator: Abstract and Section1, printed pp.1--3; inflation-spillover model, not output-synchronization identification. Primary text checked. \url{https://www.newyorkfed.org/medialibrary/media/research/staff_reports/sr1080.pdf}.
\end{itemize}
The reference models inflation spillovers; the output-correlation network intervention is a new specification.
\begin{enumerate}\raggedright
\item\hypertarget{definitions:source:30007519:1}{} Julian di Giovanni; \c{S}ebnem Kalemli-\"{O}zcan; Alvaro Silva; Muhammed A. Y\i{}ld\i{}r\i{}m. 2023. \emph{Pandemic-Era Inflation Drivers and Global Spillovers}. Abstract and Section1, printed pp.1--3; inflation-spillover model, not output-synchronization identification.
\url{https://www.newyorkfed.org/medialibrary/media/research/staff_reports/sr1080.pdf}.
DOI: \url{https://doi.org/10.59576/sr.1080}.
\end{enumerate}

\subsection*{Integrated refinement: ECON-083}
{\sffamily\small\color{muted}Multinational firms and cross-border shock transmission\par}
This refinement adopts the additional source conventions
(page~\pageref{audited:conventions}), subject to its local restrictions.
\textbf{Ownership-mediated multinational transmission.}
Distinguish multinational ownership links $G^O$ from supplier links $G^S$.
For a parent-country shock $\epsilon$ and a fixed host-country firm population,
identify host impulse responses
\[
\Psi_r^Y(h)=\left.\frac{\partial}{\partial\epsilon}
                \sum_{j\in r}Y_{j,t+h}(\epsilon)\right|_{\epsilon=0},
\qquad Y\in\{Q,L,B\},
\]
for output, employment and financing with their own units. State shock timing,
support, normalization and regularity, or use finite shock contrasts. Separate
ownership financing and knowledge channels from ordinary trade links, local
demand and common shocks; allow domestic-firm spillovers and endogenous entry.
Link-specific removal experiments require explicit replacement technologies or
financing and a common initial state, as in the existing transmission exercise.
\subsection*{Supplied source audit for ECON-083}
Local [S1], [S2], etc. in this audit and the references immediately below
belong to ECON-083, independently of the original entry's reference numbering.
[S1]'s open-questions chapter and [S2] explicitly support multinational transmission through international networks. The impulse-response and channel definitions are editorial.
\subsection*{References for integrated refinement ECON-083}
{\footnotesize\raggedright
\par\noindent\textbf{[S1]} Stefania Garetto, Nina Pavcnik, Natalia Ramondo, Vanessa Alviarez, Jingting Fan, Nitya Pandalai-Nayar, Nicola Limodio, Isabela Manelici, Nicolas Morales, Evangelina Dardati, Ezequiel Garcia-Lembergman, Grace Weishi Gu, Galina Hale, David Hémous, Ralf Martin, Farid Farrokhi, Heitor S. Pellegrina, Pierre-Louis Vézina, Laura Boudreau, and Jose P. Vasquez (2025). \emph{Foreign Direct Investment and Development}. VoxDevLit 13(1), February 2025.\par \textit{Verified locator / source role:} Conclusion/evidence-gap chapter at the linked primary review URL, inspected for the agenda. See the audit conclusion for distinctions between direct agenda support and editorial extensions.\par \url{https://voxdev.org/voxdevlit/foreign-direct-investment/fdi-advances-open-questions}\par\smallskip
\par\noindent\textbf{[S2]} Oliver Hanney / VoxDev (living compilation). \emph{Evidence gaps: Key unanswered questions in development economics}. Accessed 8 October 2026. The page currently covers 20 topics; the ``16-topics URL is a historical slug.\par \textit{Verified locator / source role:} Topic heading: Foreign Direct Investment and Development. Supports the broad research agenda; this is not a verbatim quotation or an attribution of the displayed mathematical target.\par \url{https://voxdev.org/topic/evidence-gaps-key-unanswered-questions-across-16-topics-development-economics}\par\smallskip
}

\subsection*{Common conventions: Source mathematical conventions}

These conventions apply to each entry unless explicitly replaced.

\textbf{Models and identification.} A primitive model \(m\) specifies all probability laws, feasible choices, information, equilibrium and measurement rules used by its target. The admissible class \(\mathcal M\) is fixed independently of outcomes. If \(P_O(m)\) is its observable probability law and \(\tau(m)\) the target, its identified set at observable law \(P\) is
\[
 \mathcal I_\tau(P)=\{\tau(m):m\in\mathcal M,\ P_O(m)=P\}.
\]
Point identification means that this set is a singleton. An empty set signals incompatibility of data and assumptions. Coordinate bounds need not describe the joint set. Bounds are sharp only if every retained target is attainable or arbitrarily approachable by compatible models. Finite bounds require explicit restrictions on unobserved quantities; unrestricted confounding, hidden wealth, extrapolation or arbitrary equilibrium selection can yield uninformative sets.

\textbf{Causal interventions.} A potential outcome \(Y_i(p)\) is the outcome under a complete policy regime \(p\), including timing, financing, information and market spillovers. The target population and units are fixed before treatment. With interference, use the complete assignment vector or a justified exposure mapping; individual no-interference notation alone cannot identify an economy-wide intervention. A difference of mean potential outcomes does not require the cross-world joint distribution. Conditioning on post-treatment migration, survival, enrollment or employment changes the estimand and needs separate justification.

\textbf{Statistical statements.} An estimator, confidence set, forecast or test specifies a sampling law, model class and loss/coverage criterion. Uniform \((1-\alpha)\)-coverage means \(\inf_{m\in\mathcal M}\Pr_m\{\tau(m)\in C_n\}\geq1-\alpha\), or an explicitly asymptotic version. The whole parameter space trivially has coverage, so informativeness requires a stated width, risk or power objective. A forecasting improvement is relative to a named benchmark, information set and evaluation population.

\textbf{Equilibrium and optimization.} An equilibrium comprises feasible plans, optimality, beliefs where needed, budget constraints and all market/resource clearing. The primitives or a stated benchmark define each correspondence \(\mathcal E(p;m)\). Nonemptiness is assumed or proved, never inferred just from compact policy sets. A selected-equilibrium target uses a declared selection rule; a robust target quantifies over all permitted equilibria. Use suprema or infima unless attainment is established. An \(\varepsilon\)-optimal policy is feasible and within \(\varepsilon\) of the finite optimum value. Report an empty feasible set separately.

\textbf{Welfare and accounting.} Normative utility, interpersonal normalization, social weights, numeraire, discounting and the use of public receipts are primitives. Behavioral utility need not coincide with the welfare criterion. Household welfare includes ownership income and rents once; adding profits again to compensating variation generally double-counts them. A monetary equivalent is defined by an explicit transfer and price convention, with monotonicity and range conditions. Average effects, mechanism contributions, transfers and welfare losses are different targets. Infinite sums require convergence and dynamic feasible plans require terminal or no-Ponzi conditions.

\textbf{Source versions.} A working-paper date, revision date and journal publication year are distinguished. A DOI for a journal version is not automatically the DOI of an earlier manuscript. Locators refer to the stated version and printed pages where specified. A metadata-only check does not verify a claim about a full-text conclusion. Every entry's source subsection narrows the provenance claim accordingly.

\subsection*{Common conventions: Additional-source status and mathematical conventions}

\label{audited:conventions}

Of the 100 supplied ECON entries, 65 distinct problems appear here; 32 close
matches refine earlier entries, and three duplicate questions map to existing
entries. Appendix~\ref{app:audited-crosswalk} lists every destination.
The attachment's P/R/S/U distinctions and source audits are retained. Its
precise mathematical formalizations are editorial unless the individual audit
says otherwise. Candidate subsidy targets ECON-004--006 retain their U flags;
the resolved approval-core question ECON-007 updates OP-0202 above.
The following supplied conventions govern both this part and the attributed
ECON refinements elsewhere in the catalogue, subject to local restrictions.

\textbf{Parameterized questions.} A problem family lists inputs that must be fixed before an instance is evaluated: population, horizon, policies, information, data law, model class and welfare weights. These are required choices, not quantities the citations are assumed to specify. Undefined applications should not be treated as identified estimands.

\textbf{Indices and integrability.} Sets of agents, firms and regions are finite unless a population measure is explicitly used. \(i,j\) index units, \(r\) regions, \(t\) dates and \(h\) horizons. \(\mathbb E\) and \(\Pr\) refer to the declared population and law; \(\mathbf1\{A\}\) is an indicator. Every displayed expectation and discounted sum needs existence in the stated domain. Infinite horizons use a discount in \((0,1)\) and a convergence condition; finite horizons may use discount one.

\textbf{Optimization and existence.} A displayed \(\max\), \(\min\), or \(\arg\max\) is conditional on attainment. For finite feasible menus this is immediate; general spaces need continuity and compactness/coercivity or another existence argument. Without attainment, the target is the supremum/infimum and arbitrarily near-optimal feasible choices. \(\inf\varnothing=+\infty\). Empty feasibility and an unidentified target are different issues.

\textbf{Potential outcomes.} \(Y_i(z)\) is the outcome under a specified intervention, not the observed conditional mean among people choosing \(z\). In binary comparisons \(\Delta Y_i=Y_i(1)-Y_i(0)\). Specify assignment, treatment versions, compliance, information and observation horizon. Interference is allowed under a full policy vector; compressed exposure notation needs a justified mapping. No interference, consistency and overlap are substantive assumptions, not automatic consequences of notation.

\textbf{Populations and observation.} Fix baseline eligibility and population weights. Conditioning on post-policy employment, survival, relocation or program uptake can change the population being compared. Death is not ordinary loss to follow-up; outcomes truncated by death need a composite convention, joint survival target, or explicitly defined survivor stratum. Fertility-changing interventions need a population functional rather than an unexplained fixed descendant cohort. Measurement and missingness models must be supplied.

\textbf{Identification.} A structural model \(m\in\mathcal M\) implies observable law \(P_m\) and target \(\theta(m)\). Identification means
\[
 P_{m_1}=P_{m_2}\ \Longrightarrow\ \theta(m_1)=\theta(m_2)
 \quad\text{for all }m_1,m_2\in\mathcal M .
\]
Without identification, the target is
\[
 \Theta(P;\mathcal M)=\{\theta(m):m\in\mathcal M,\ P_m=P\}.
\]
Writing \(\theta(P)\) before establishing identification can hide incompatible latent models. Confidence sets additionally need a sampling law, confidence level, asymptotic sequence and uniformity class. Point identification, coverage and informative precision are separate properties.

\textbf{Mediation.} Total effects of randomized assignment are distinct from cross-world natural indirect effects. Belief/effort-channel effects require a meaningful mediator intervention and additional measurement, exclusion or mediation assumptions. Randomization of the initial policy alone does not supply those assumptions. Report bounds or interventional alternatives when the stronger target is unsupported.

\textbf{Equilibrium.} A counterfactual \(W(z)\), \(p(z)\), or \(\mu_t^z\) requires individual optimization, feasibility, government/resource budgets, market clearing and state-transition laws. State equilibrium existence and selection, or report ranges over compatible equilibria. Initial-state integration includes subsequent endogenous transitions; current-state integration must use the policy-dependent distribution. Reduced-form invariance after policy is not assumed.

\textbf{Welfare accounts.} Scores, utility, health, dollars and ecological quantities require explicit conversion before addition. Fees, taxes, subsidies, wages and extortion payments are transfers between parties; distributional weights can make them welfare relevant, but their face value is not automatically a resource loss. State ownership, geographic boundaries, foreign-income weights and fiscal finance. Do not add profits to household welfare already including dividends, or count land capitalization and the underlying benefit twice. A benefit-cost expression must distinguish gross benefits from net welfare and subtract every real cost exactly once.

\textbf{Derivatives and risk.} Log derivatives require positive arguments; local responses require admissible perturbations and specified equilibrium branches. Differentiation under expectations needs regularity/domination, and boundaries may allow only one-sided derivatives. Risk uses a specified law; ambiguity uses a declared nonempty law set on a common history space. Adaptive policies must be nonanticipative. A robust criterion is an editorial choice unless attributed explicitly.

\textbf{Scope overlaps.} Allocation, collective choice, contracts, household bargaining and coordinated policy overlap economics with optimization and game theory. They are retained because they appeared in the original catalogue; no claim of a strictly game-theory-free selection is made.
% END ECONOMICS STATEMENT
\end{document}
